\bookchapter{Глава 2: Наведи порядок в своей цифровой комнате — управление файлами}

\chapterintro{Со временем телефон захламляется, и хорошие привычки начинаются с управления файлами. Так же как ваша комната нуждается в регулярной уборке, ваше цифровое пространство тоже требует внимания. Хорошо организованный телефон не только облегчает поиск, но и работает более плавно.}

\booksection{I. Трёхуровневая система хранения: ваш цифровой склад}

\begin{enumerate}
  \item \textbf{Встроенная память}: основная зона, высокая скорость, здесь находятся система и приложения. Это самое ценное пространство вашего телефона, расходуйте его экономно.
  \item \textbf{SD-карта (если поддерживается)}: расширенный шкаф, скорость ниже, подходит для хранения фильмов, фото и других крупных файлов. Внимание: многие флагманские модели сейчас не поддерживают SD-карты.
  \item \textbf{Облачное хранилище}: сейф + синхронизируемый сетевой диск. Защита от потери и синхронизация между устройствами. iCloud, Baidu Netdisk, OneDrive — популярные варианты.
\end{enumerate}

\booksection{II. Знакомство с «Файловым менеджером»}

На Android обычно есть приложение [Файлы] или [Файловый менеджер], на iOS — встроенное приложение [Файлы]. Это ваш «Проводник» для просмотра всех файлов телефона.

\booksubsection{Значение распространённых папок:}

\begin{itemize}
  \item \textbf{DCIM}: фото и видео, снятые камерой.
  \item \textbf{Pictures}: другие изображения — скриншоты, сохранённые картинки.
  \item \textbf{Download}: файлы, загруженные через браузер.
  \item \textbf{Music}: музыкальные файлы.
  \item \textbf{Movies}: видеофайлы.
  \item \textbf{Documents}: файлы документов.
  \item \textbf{Android}: системная папка Android, внутри — данные приложений. Обычно не удаляйте вручную.
\end{itemize}

\booksection{III. Структура файловой системы: углублённое понимание}

Android основан на Linux, файловая система имеет определённую иерархическую структуру. Хотя обычному пользователю не нужно глубоко вникать, знание некоторых базовых концепций полезно.

\booksubsection{Разделы памяти:}

\begin{itemize}
  \item \textbf{/system}: системный раздел, содержит файлы операционной системы. Обычный пользователь не имеет доступа, требуются права root.
  \item \textbf{/data}: раздел данных, содержит пользовательские данные и данные приложений. Это пространство, которое вы используете ежедневно.
  \item \textbf{/sdcard или /storage/emulated/0}: доступное пользователю пространство памяти, здесь находятся ваши фото и загруженные файлы.
\end{itemize}

\booksubsection{Хранение данных приложений:}

\begin{itemize}
  \item \textbf{Внутреннее хранилище}: приватное пространство приложения, недоступное другим приложениям. Удаляется при удалении приложения.
  \item \textbf{Внешнее хранилище}: файлы, сохранённые приложением в доступной пользователю области, например фото и загрузки. Не обязательно удаляются при удалении приложения.
\end{itemize}

\booksection{IV. Резервное копирование и восстановление данных: страховка для ваших данных}

Данные бесценны, резервное копирование — обязательная привычка.

\booksubsection{Стратегии резервного копирования:}

\begin{enumerate}
  \item \textbf{Облачное резервное копирование}: автоматическое, удобное, но требует сеть и облачное пространство. iCloud, Huawei Cloud, Xiaomi Cloud и другие поддерживают.
  \item \textbf{Локальное резервное копирование}: подключите телефон к компьютеру кабелем и используйте официальное ПО для резервного копирования. Быстро, много места, но требует ручных действий.
  \item \textbf{Двойное резервное копирование}: облачное + локальное — самое надёжное.
\end{enumerate}

\booksubsection{Что копировать?}

\begin{itemize}
  \item \textbf{Фото и видео}: занимают больше всего места и самые ценные.
  \item \textbf{Контакты}: удобнее всего хранить в облачной учётной записи.
  \item \textbf{История чатов WeChat}: используйте встроенную функцию резервного копирования WeChat.
  \item \textbf{Важные документы}: храните в облачном диске или на компьютере.
\end{itemize}

\booksection{V. Безопасность файлов: защитите вашу приватность}

\begin{itemize}
  \item \textbf{Шифрование чувствительных файлов}: многие телефоны имеют функцию «Безопасная папка» или «Приватное пространство», куда можно поместить чувствительные файлы, доступ к которым требует пароль или отпечаток пальца.
  \item \textbf{Удаление не равно уничтожению}: обычное удаление лишь помечает пространство как доступное для перезаписи, данные всё ещё там. Для полного удаления нужны специальные инструменты безопасного стирания или многократная перезапись.
  \item \textbf{Безопасность в облаке}: выбирайте надёжных облачных провайдеров, включите двухфакторную аутентификацию, регулярно проверяйте устройства входа.
\end{itemize}

\booksection{VI. Не хватает места? Устроим генеральную уборку}

\begin{itemize}
  \item \textbf{Очистка кэша}: временная память приложений, очистка освобождает место. Можно очистить в Настройки > Память.
  \item \textbf{Поиск больших файлов}: используйте файловый менеджер для сортировки по размеру и найдите «виновников», занимающих место.
  \item \textbf{Удаление остатков}: после удаления приложений иногда остаются мусорные файлы — удалите вручную.
  \item \textbf{Очистка галереи}: удалите дубликаты фото, размытые снимки и ненужные скриншоты.
  \item \textbf{Очистка WeChat}: WeChat — главный пожиратель места, регулярная очистка истории чатов и кэша очень необходима.
\end{itemize}

\booksection{VII. Полный разбор распространённых форматов файлов}

Различные расширения файлов в телефоне сбивают с толку, но разобравшись в них, вы сможете лучше управлять файлами:

\begin{center}
\begin{tabularx}{\linewidth}{|X|X|X|X|}
\hline
\textbf{Формат} & \textbf{Тип} & \textbf{Особенности} & \textbf{Рекомендуемое приложение} \tabularnewline
\hline
JPG/JPEG & Изображение & Сжатие с потерями, малый размер, лучшая совместимость. Формат по умолчанию для фото с телефона. & Галерея / любой просмотрщик \tabularnewline
PNG & Изображение & Сжатие без потерь, поддержка прозрачного фона. Часто используется для скриншотов и иконок. Размер больше JPG. & Галерея / любой просмотрщик \tabularnewline
MP3 & Аудио & Сжатие с потерями, малый размер, отличная совместимость. Самый универсальный аудиоформат. & Любой музыкальный плеер \tabularnewline
FLAC & Аудио & Сжатие без потерь, качество как у CD, размер в 3-5 раз больше MP3. Выбор аудиофилов. & Музыкальный плеер / проф. приложение \tabularnewline
WAV & Аудио & Несжатое исходное аудио, лучшее качество, но огромный размер. Часто в профессиональной записи. & Музыкальный плеер \tabularnewline
M4A & Аудио & Кодек AAC, при том же битрейте качество лучше MP3. Часто в экосистеме Apple. & Apple Music / музыкальный плеер \tabularnewline
MP4 & Видео & Самый универсальный видеоформат, отличная совместимость. Формат по умолчанию для видео с телефона. & Галерея / любой видеоплеер \tabularnewline
MKV & Видео & Открытый формат, поддержка нескольких аудиодорожек и субтитров. Часто для фильмов. & VLC / nPlayer \tabularnewline
DOC/DOCX & Документ & Документы Word. DOCX — сжатый формат, меньший размер. DOC — старый формат. & WPS / Office \tabularnewline
PPT/PPTX & Презентация & Презентации PowerPoint. PPTX — сжатый формат. & WPS / Office \tabularnewline
XLS/XLSX & Таблица & Электронные таблицы Excel. XLSX — сжатый формат. & WPS / Office \tabularnewline
PDF & Документ & Кроссплатформенный формат, фиксированная вёрстка, не редактируется. Для договоров, статей, электронных книг. & Adobe Reader / WPS \tabularnewline
ZIP & Архив & Самый универсальный формат сжатия, хорошая кроссплатформенная совместимость. & Системный / архиватор \tabularnewline
RAR & Архив & Степень сжатия выше, чем у ZIP, но требуется специальное ПО для распаковки. & Архиватор \tabularnewline
ISO & Образ & Файл образа диска, полное копирование содержимого диска. Часто для установочных пакетов системы. & Архиватор / виртуальный привод \tabularnewline
HTML & Веб-страница & Файл веб-страницы, открывается браузером. Сохранённые веб-страницы часто в этом формате. & Браузер \tabularnewline
\hline
\end{tabularx}
\end{center}

\begin{tcbpractice}{Упражнение: наведите порядок в папке «Загрузки»}
Откройте файловый менеджер > зайдите в папку «Загрузки». Посмотрите, что там есть? Удалите ненужное, полезное рассортируйте по категориям (например, создайте папки «Рабочие материалы», «Стикеры»). Заодно проверьте использование памяти и узнайте, какие типы файлов занимают больше всего места.

\end{tcbpractice}

\booksection{VIII. Управление изображениями и файловой системой: наведём порядок в активах}

В предыдущем разделе мы научились «узнавать папки», но одного знания недостаточно. В этом разделе мы расскажем,\textbf{как активно создать систему управления, которой можно следовать в долгосрочной перспективе}, чтобы ваши файлы всегда можно было найти, они не терялись и не были доступны посторонним.

\booksubsection{Классификация и соглашения об именовании файлов}

Корень хаоса — отсутствие правил. Создать правила на самом деле просто, главное —\textbf{придерживаться их}.

\begin{itemize}
  \item \textbf{Именование «Время + Событие»}: в названия папок включайте временную метку и описание события для удобного последующего поиска. Например:\texttt{202606\_Чеки\_путешествия},\texttt{2025\_Годовой\_отчёт},\texttt{2024\_День\_рождения\_ребёнка}. Так вы видите и хронологию, и сразу понимаете содержание.
  \item \textbf{Трёхуровневая система классификации}:
  \begin{itemize}
    \item \textbf{Уровень 1}: по типу файла (документы / изображения / видео / аудио)
    \item \textbf{Уровень 2}: по сценарию использования (работа / жизнь / учёба)
    \item \textbf{Уровень 3}: по конкретному событию или проекту (например, «2026Q2\_Проект\_А»)
  \end{itemize}
  \item \textbf{Ежемесячный план уборки}: каждый месяц в фиксированный день (рекомендуется в день зарплаты или 1-го числа) проводите централизованную уборку каталога «Загрузки»:\textbf{архивируйте полезное, удаляйте бесполезное, сохраняйте ценное}. 12 маленьких уборок в год гораздо легче, чем одна генеральная уборка раз в год.
\end{itemize}

\begin{tcbpractice}{Список задач на этот месяц}
Выберите свободные выходные и действуйте в следующем порядке:

\begin{enumerate}
  \item Создайте в файловом менеджере рабочую папку на этот месяц в формате\texttt{2026Q2\_Проект\_А};
  \item Просмотрите все файлы в каталоге «Загрузки» за эту неделю, оставьте полезное, заархивируйте, удалите ненужное;
  \item Создайте годовые папки для семейных фотографий:\texttt{2025\_Семья},\texttt{2026\_Семья}.
\end{enumerate}

\end{tcbpractice}

\booksubsection{Стратегия обработки избыточных файлов}

Со временем телефон «толстеет», и чаще всего виноваты избыточные файлы.

\begin{itemize}
  \item \textbf{Автоочистка за 72 часа}: временные файлы, полученные в социальных приложениях, таких как WeChat и QQ (стикеры, изображения из чатов, короткие видео), в большинстве своём после одного просмотра бесполезны. В [WeChat] > [Настройки] > [Общие] > [Память] включите «Автоочистку» и установите срок хранения кэша 72 часа. Телефон сразу станет заметно легче, при этом важные записи не потеряются.
  \item \textbf{Целевая очистка памяти}: Android [Настройки] > [Память] > [Очистка и ускорение], iOS [Настройки] > [Общие] > [Хранилище iPhone]. Система покажет категории: «Кэш / Временные файлы / Удаляемые приложения / Крупные файлы». Сосредоточьтесь на очистке: (1) кэш-файлов; (2) просроченных установочных пакетов (старые версии, оставшиеся после обновления приложений); (3) дубликатов файлов (фото с одинаковым именем и размером).
  \item \textbf{Правила автоматической классификации}: большинство фирменных приложений галереи поддерживают умную классификацию: люди, пейзажи, документы, скриншоты. После включения новые групповые фото автоматически попадут в «Люди > Иван», отсканированные договоры — в «Документы».\textbf{В сочетании с ежемесячной уборкой это значительно снижает вероятность «не могу найти».}
\end{itemize}

\booksubsection{Правила безопасной обработки изображений}

Изображения — самый «чувствительный» тип данных в телефоне: один случайный снимок может раскрыть ваше местоположение, устройство и время. Мы рекомендуем\textbf{создать рабочий процесс «двойной защиты»}:

\begin{tcbwarn}{Стандартный процесс отправки фото документов / чувствительных скриншотов}
\begin{enumerate}
  \item \textbf{Закрашивание видимой информации}: с помощью встроенных или сторонних инструментов (например, Meitu Xiuxiu, Snapseed)\textbf{замажьте средние цифры номера паспорта, номера автомобиля, адрес, имя, сумму}и другие ключевые поля;
  \item \textbf{Удаление метаданных}: перед отправкой в галерее нажмите «Поделиться» или «Экспорт», система спросит, сохранять ли данные о местоположении — выберите\textbf{«Не сохранять»}. Тогда модель устройства, GPS-координаты и время съёмки не уйдут вместе с фото;
  \item \textbf{Временный водяной знак}(опционально): нанесите на чувствительный скриншот текстовый водяной знак «Только для оформления ХХ», чтобы избежать злоупотребления.
\end{enumerate}

\end{tcbwarn}

Распространённые методы очистки метаданных (EXIF) (базовые операции, для глубокой защиты см. главу 6 «Обезличивание фотографий»):

\begin{itemize}
  \item \textbf{Системные средства}: iOS [Настройки] > [Приватность и безопасность] > [Камера] — можно отключить «Точное местоположение»; при отправке выберите «Только текущее фото» —\textbf{автоматически удалит большую часть EXIF}.
  \item \textbf{Android}: большинство фирменных приложений галереи при «Поделиться» спрашивают «Удалить информацию о местоположении» — отметьте галочкой.
  \item \textbf{Профессиональные инструменты}: ExifTool (ПК), Metapho (iOS), Photo Metadata Remover (Android) могут полностью удалить все метаданные.
\end{itemize}

Для чувствительных изображений, требующих\textbf{долгосрочного хранения}(паспорт, банковская карта, скан домовой книги, загранпаспорт), рекомендуется хранить в\textbf{зашифрованной галерее}:

\begin{itemize}
  \item iPhone: [Фото] > выберите фото > Скрыть (доступно только после проверки Face ID);
  \item Huawei / Xiaomi: [Галерея] > долгое нажатие на фото > «Переместить в сейф»;
  \item Сторонние приложения для шифрования:\textbf{360 Security Album},\textbf{Tencent Gallery Manager (Private Gallery)},\textbf{KeepSafe}и др. — отдельная защита паролем, лучшая скрытность.
\end{itemize}
